\documentclass[aspectratio=169,10pt]{beamer}
\usetheme[numbering=fraction,progressbar=frametitle]{metropolis}
\usepackage{fontspec}
\setsansfont{Latin Modern Sans}
\usepackage{amsmath,amssymb,mathtools,graphicx,booktabs}
\usepackage{polyglossia}
\setdefaultlanguage{french}
\setbeamertemplate{navigation symbols}{}
\setbeamerfont{frametitle}{size=\large}
\setbeamerfont{normal text}{size=\normalsize}
\setlength{\parskip}{.45em}
\hypersetup{pdfauthor={},pdfsubject={Lecture slides},pdfkeywords={},pdftitle={Droites et cercles}}
\providecommand{\N}{\mathbb{N}}
\providecommand{\Z}{\mathbb{Z}}
\providecommand{\D}{\mathbb{D}}
\providecommand{\Q}{\mathbb{Q}}
\providecommand{\R}{\mathbb{R}}
\providecommand{\C}{\mathbb{C}}
\providecommand{\vect}[1]{\overrightarrow{#1}}
\providecommand{\vi}{\vec{\imath}}
\providecommand{\vj}{\vec{\jmath}}
\providecommand{\vk}{\vec{k}}
\providecommand{\norm}[1]{\left\lVert #1\right\rVert}
\providecommand{\vc}[2]{\begin{pmatrix}#1\\#2\end{pmatrix}}
\providecommand{\deter}[2]{\det\left(#1,#2\right)}
\providecommand{\abs}[1]{\left\lvert #1\right\rvert}
\providecommand{\dd}{\,\mathrm{d}}
\providecommand{\e}{\mathrm{e}}
\providecommand{\Card}{\operatorname{Card}}
\providecommand{\E}{\operatorname{E}}
\providecommand{\V}{\operatorname{V}}
\providecommand{\degre}{^{\circ}}
\providecommand{\un}[1]{\,\mathrm{#1}}
\providecommand{\sci}[2]{#1\times10^{#2}}
\providecommand{\pgcd}{\operatorname{PGCD}}
\providecommand{\ppcm}{\operatorname{PPCM}}
\title{Droites et cercles}
\subtitle{Cours et exemples}
\date{}
\author{}
\begin{document}
\begin{frame}[plain]
\vfill
{\usebeamerfont{title}\inserttitle\par}
\medskip
{\usebeamerfont{subtitle}\insertsubtitle\par}
\medskip
{\usebeamercolor[fg]{progress bar}\rule{\textwidth}{.5pt}}
\vfill
\end{frame}
\begin{frame}{}
Marquer un point \(\Omega\) puis placer le plus grand nombre possible de points à
exactement \(4\) cm de \(\Omega\). Décrire l'ensemble obtenu. Dans un repère où
\(\Omega(a;b)\), traduire la contrainte de distance par une équation.
\end{frame}
\begin{frame}{Objectifs}
\begin{itemize}

  \item Écrire les équations d'une droite sous toutes ses formes.
  
\item Écrire et reconnaître l'équation d'un cercle.
  
\item Calculer la distance d'un point à une droite.
  
\item Étudier la position relative d'une droite et d'un cercle.

\end{itemize}
\end{frame}
\begin{frame}{Équations d'une droite}
Dans un repère :

\begin{itemize}

  \item \textbf{Cartésienne} : \(ax + by + c = 0\), avec \((a,b)\ne(0,0)\).
  
\item \textbf{Réduite} : \(y = mx + p\) (\(m\) = coefficient directeur).
  
\item \textbf{Paramétrique} :
        \(\begin{cases}x = x_0 + t\,d_x\\ y = y_0 + t\,d_y\end{cases}\),
        \(t\in\mathbb{R}\), \((d_x,d_y)\ne(0,0)\) vecteur directeur.

\end{itemize}
\end{frame}
\begin{frame}{Cercles et distance}
Cercle de centre \(\Omega(a,b)\) et rayon \(r>0\) :
\[(x-a)^2 + (y-b)^2 = r^2.\]
\end{frame}
\begin{frame}{Cercles et distance}
Distance du point \(A(x_0,y_0)\) à la droite \(ax+by+c=0\) :
\[d\bigl(A,(\Delta)\bigr) = \dfrac{\left|ax_0+by_0+c\right|}{\sqrt{a^2+b^2}}.\]
\end{frame}
\begin{frame}{Cercles et distance}
Pour le cercle de centre \(\Omega\), rayon \(r\), et la droite \((\Delta)\) :

\begin{itemize}

  \item \(d(\Omega,\Delta) > r\) : extérieure (0 point commun).
  
\item \(d(\Omega,\Delta) = r\) : tangente (1 point commun).
  
\item \(d(\Omega,\Delta) < r\) : sécante (2 points communs).

\end{itemize}
\end{frame}
\begin{frame}{Droites : propriétés et méthodes}
Pour une droite \(y=mx+p\) :

\begin{itemize}

  \item \(m>0\) : droite croissante (monte de gauche à droite).
  
\item \(m<0\) : droite décroissante.
  
\item \(m=0\) : droite horizontale.
  
\item Droite verticale : équation \(x=k\) (pas de coefficient directeur).

\end{itemize}
\end{frame}
\begin{frame}{Droites : propriétés et méthodes}
\begin{enumerate}
\setcounter{enumi}{0}
  \item Par deux points \(A(x_A,y_A)\) et \(B(x_B,y_B)\) :
        \(m = \dfrac{y_B-y_A}{x_B-x_A}\) (si \(x_A\ne x_B\)),
        puis \(p = y_A - mx_A\).
  
\item \textbf{Par un point et une pente :}
        \(y - y_A = m(x-x_A)\).
  
\item Parallèle à \(y=mx+p\) : même pente \(m\), chercher \(p'\)
        via le point donné.
  
\item Perpendiculaire à \(y=mx+p\) : si \(m\ne0\), pente \(m'=-1/m\).
        Si \(m=0\), la perpendiculaire passant par \(A\) est la verticale \(x=x_A\).

\end{enumerate}
\end{frame}
\begin{frame}{Droites : propriétés et méthodes}
Deux droites \(d_1\) : \(y=m_1x+p_1\) et \(d_2\) : \(y=m_2x+p_2\) :

\begin{itemize}

  \item \(d_1 \parallel d_2 \Leftrightarrow m_1=m_2\).
  
\item \(d_1 \perp d_2 \Leftrightarrow m_1\times m_2 = -1\).

\end{itemize}
\end{frame}
\begin{frame}{Droites : propriétés et méthodes}
La médiatrice de \([AB]\) (\(A(1,3)\), \(B(5,1)\)) passe par le milieu
\(I(3,2)\) et est perpendiculaire à \(\overrightarrow{AB}=(4,-2)\),
donc son coefficient directeur est \(m=2\).
Équation : \(y-2 = 2(x-3)\), soit \(y=2x-4\).
\end{frame}
\begin{frame}{Cercles : propriétés et méthodes}
L'équation \(x^2+y^2+Dx+Ey+F=0\) représente un cercle si et seulement si
\(D^2+E^2-4F > 0\).
On le met sous forme standard en complétant le carré :
\[\left(x+\dfrac{D}{2}\right)^2+\left(y+\dfrac{E}{2}\right)^2 = \dfrac{D^2+E^2-4F}{4}.\]
Centre : \(\Omega\!\left(-\dfrac{D}{2},\,-\dfrac{E}{2}\right)\) ;
rayon : \(r=\dfrac{\sqrt{D^2+E^2-4F}}{2}\).
\end{frame}
\begin{frame}{Cercles : propriétés et méthodes}
La tangente au cercle \(\mathcal{C}\) de centre \(\Omega\) en un point \(A\) de \(\mathcal{C}\) :

\begin{itemize}

  \item Est perpendiculaire au rayon \(\Omega A\).
  
\item Passe par \(A\).

\end{itemize}

Pour \(\mathcal{C}\) : \(x^2+y^2=r^2\) et \(A(x_0,y_0)\in\mathcal{C}\) :
équation de la tangente : \(x_0 x + y_0 y = r^2\).
\end{frame}
\begin{frame}{Cercles : propriétés et méthodes}
Deux cercles de centres \(\Omega_1\), \(\Omega_2\) et de rayons \(r_1\), \(r_2\)
avec \(d = \Omega_1\Omega_2\) :

\begin{itemize}

  \item \(d > r_1+r_2\) : extérieurs l'un à l'autre (0 point commun).
  
\item \(d = r_1+r_2\) : tangents extérieurement (1 point commun).
  
\item \(|r_1-r_2| < d < r_1+r_2\) : sécants (2 points communs).
  
\item \(d = |r_1-r_2|>0\) : tangents intérieurement (1 point commun).
  
\item \(d=0\) et \(r_1=r_2\) : cercles confondus (une infinité de points communs).
  
\item \(d < |r_1-r_2|\) : l'un à l'intérieur de l'autre (0 point commun).

\end{itemize}
\end{frame}
\end{document}
